\Needspace{18\baselineskip}
\section{Typed boxes, deep occurrences, and contexts}
\Needspace{24\baselineskip}
\subsection{Relabelling without changing incidence}
A signature-preserving map $f$ of operation alphabets relabels a net by replacing
each edge label $o$ with $f(o)$. All carriers, enumerations, vertex labels, and
incidence data are retained. Hence their invariant proofs are retained too.
For typing, compose the old typing equality with the inverse of the stipulated
signature equality for $f(o)$. This yields the new signature list at that edge.
\lean{Graph/Hierarchy}{17}{36}{Net.relabel}

\Needspace{24\baselineskip}
\subsection{The linear wire and label discipline}
\begin{definition}[Wire types, operation arities, and holes]
A wire type is either a term wire $*$ or a thunk wire $T^n(*)$, carrying a natural
number $n$ of bound variables. An operation specification assigns each operation
an argument count $a(o)$ and an ordered list of binder counts $b(o)$. A hole label
contains a natural-number name together with arbitrary input and output wire lists.
Equality of hole labels includes these lists as well as the name.
\end{definition}
\lean{Graph/Hierarchy}{44}{47}{Linear.Wire}
\lean{Graph/Hierarchy}{51}{53}{Linear.Operations}
\lean{Graph/Hierarchy}{56}{60}{Linear.Hole}
Primitive edge labels are operations, variables, focus labels from the abstract
three-mode type, or typed holes. This introduces concrete focus \emph{labels};
it does not yet define the UAM transition relation.
\lean{Graph/Hierarchy}{62}{66}{Linear.Primitive}
Their signature is as follows. Every non-hole primitive has one term input.
Operation $o$ has $a(o)$ term outputs followed by the thunk outputs specified in
$b(o)$. A variable or focus has one term output. A hole uses precisely the input
and output lists stored in its label. Ordering is part of each of these definitions.
\lean{Graph/Hierarchy}{68}{76}{Linear.primitiveSignature}

\begin{definition}[Alphabets and typed boxes]
An alphabet is a label type and a signature on it. A box over alphabet $A$ stores
numbers $n,m$ and a body net over $A$ of type
\[
 *\longrightarrow *^{n+m}.
\]
The corresponding outer box edge will have type
$T^n(*)\longrightarrow *^m$. Thus body boundary and outer boundary differ:
$n$ body outputs account for binding, and $m$ for exposed free variables.
\end{definition}
\lean{Graph/Hierarchy}{80}{82}{Linear.Alphabet}
\lean{Graph/Hierarchy}{84}{87}{Linear.Box}
All alphabet label types live in one fixed higher universe, while graph carriers
remain small types. The base alphabet contains just primitives, raised into that
universe by a universe lift; its signature reads the underlying primitive. The
successor alphabet is the disjoint union of primitives and boxes over the previous
alphabet. Its signature uses the primitive signature on the left and the stated
external box interface on the right.
\lean{Graph/Hierarchy}{89}{94}{Linear.baseAlphabet}
\lean{Graph/Hierarchy}{96}{105}{Linear.nextAlphabet}
Recursively, depth zero uses the base alphabet and depth $d+1$ uses the successor
of the depth-$d$ alphabet. A hypernet at depth bound $d$ is a net over that
alphabet. The bound limits possible nesting; it need not be the actual minimum
depth of the particular graph.
\lean{Graph/Hierarchy}{107}{109}{Linear.alphabet}
\lean{Graph/Hierarchy}{111}{112}{Linear.Hypernet}

The operation predicate selects primitive operation labels only. At depth zero
it inspects the universe-lifted primitive; at a successor depth it accepts only
the primitive-operation summand, excluding boxes and all other primitive kinds.
\lean{Graph/Hierarchy}{114}{118}{Linear.IsOperation}
Selected operation edges have exactly one source. Rewrite the incidence length
using signature typing, split on the depth and label constructors, discard labels
excluded by the predicate, and compute the operation input list $[*]$.
\lean{Graph/Hierarchy}{120}{131}{Linear.unary_operations}
It follows that two operation paths from inputs to the same vertex have identical
edge lists and input indices in every typed hypernet. Instantiate the generic
operation-path uniqueness theorem with the preceding derived unary proof. This
is the graph-theoretic uniqueness conclusion associated with the paper's Lemma~22
for this representation; neither determinism nor a locality field is assumed.
\lean{Graph/Hierarchy}{135}{140}{Linear.operation_path_unique}

\Needspace{24\baselineskip}
\subsection{Box construction and promotion of the depth bound}
The box constructor takes a body with the specified $n+m$ output interface and
forms the corolla of its box label in the successor alphabet. Corolla
well-formedness proves all outer graph conditions. The result has exactly one
outer box edge and the external type $T^n(*)\to *^m$.
\lean{Graph/Hierarchy}{143}{147}{Linear.box}
Promotion raises the available depth bound by one without changing the outer
graph. At depth zero, send a lifted primitive to the primitive summand at depth
one. At a successor depth, keep primitives and recursively promote each box
body, retaining its bound/free metadata. These maps preserve signatures because
neither primitive signatures nor external box arities change. Apply the generic
relabel construction at each recursive step.
\lean{Graph/Hierarchy}{151}{161}{Linear.raise}
Outer sources and targets are unchanged by promotion. The code records these
as heterogeneous equalities: the dependent graph types are indexed by different
depths, although after unfolding either defining case the incidence functions
are definitionally the same. Splitting on depth proves both equalities by
reflexivity.
\lean{Graph/Hierarchy}{163}{165}{Linear.raise_source}
\lean{Graph/Hierarchy}{167}{169}{Linear.raise_target}

\Needspace{24\baselineskip}
\subsection{Finite occurrence addresses at arbitrary depth}
\begin{definition}[All edge occurrences]
At depth zero, an occurrence is a shallow edge. At depth $d+1$, an occurrence is
either a shallow edge or a pair $(e,a)$ where $e$ is a box edge and $a$ an
occurrence in its body. A primitive edge has the empty type of child occurrences.
\end{definition}
\lean{Graph/Hierarchy}{173}{178}{Linear.AllEdge}
The parent edge is part of every deep address. In particular identical bodies
placed at different parent edges do not identify their child occurrences.

\begin{theorem}[Finiteness of all occurrences]
Every finite-depth hypernet has an explicit finite enumeration of its shallow
and deep edge occurrences.
\end{theorem}
\begin{proof}
Recurse on the depth bound. At zero use the graph's edge enumeration. At a
successor bound, enumerate shallow edges in the left summand. For the right
summand enumerate the dependent sum over outer edges: a primitive contributes
an empty enumeration, and a box contributes the recursively obtained enumeration
of its body. The sum and dependent-sum enumeration constructions already prove
completeness of this list. Recursion decreases the depth at every box descent.
\end{proof}
\lean{Graph/Hierarchy}{182}{190}{Linear.allEdgesFinite}
Every given occurrence belongs to this enumeration, by its completeness field.
This is the finiteness content of Appendix~A, Lemma~19, for the typed hierarchy.
The list need not be duplicate-free, since completeness suffices here.
\lean{Graph/Hierarchy}{192}{195}{Linear.every_edge_enumerated}

\Needspace{24\baselineskip}
\subsection{Hole lookup and the context condition}
A primitive hole lookup returns its stored label for a hole constructor and no
label otherwise. Full occurrence lookup uses that function on a shallow primitive,
returns no hole at a shallow box edge, and recursively looks inside a deep box
address. A purported child occurrence under a primitive has empty type and is
eliminated. Dependent transports in Lean express that the selected label really
is the box whose body is being searched.
\lean{Graph/Hierarchy}{197}{199}{Linear.primitiveHole}
\lean{Graph/Hierarchy}{202}{214}{Linear.holeAt}
\begin{definition}[Unique holes and contexts]
A hypernet has unique holes if, for each complete typed hole label $h$, any two
occurrences whose lookup returns $h$ are equal. A context is a hypernet together
with this property. This permits multiple different hole labels and also permits
no holes; it is an at-most-once condition for each label over all depths.
\end{definition}
\lean{Graph/Hierarchy}{218}{221}{Linear.UniqueHoles}
\lean{Graph/Hierarchy}{223}{224}{Linear.Context}
\begin{lemma}[Boxing preserves unique holes]
If a body has unique holes, its one-edge box has unique holes.
\end{lemma}
\begin{proof}
A shallow occurrence in the outer box is not a hole, so either proposed hole
occurrence must be deep. The outer graph has only one edge; therefore their
parent indices agree. Their lookups are the body lookups, whose occurrences
agree by the body's uniqueness assumption. Tagging these equal child addresses
with their common parent gives equality of the original occurrences.
\end{proof}
\lean{Graph/Hierarchy}{228}{249}{Linear.box_unique_holes}
Boxing a context now means boxing its underlying graph and supplying this lemma
as the context certificate. Shallow plugging at any fixed hierarchy depth is
just the generic typed net plugging construction with that depth's alphabet;
its interface constraints are checked by the type of its replacement argument.
\lean{Graph/Hierarchy}{251}{254}{Linear.boxContext}
\lean{Graph/Hierarchy}{257}{261}{Linear.plug}
This last operation returns a hypernet. The later context-calculus development
now proves that its result is a context under explicit freshness conditions
(Section~\ref{sec:context-closure}), including a sharp necessary-and-sufficient
condition for shallow plugging. The typed hierarchy has not yet been proved
adequate for the paper's original set-based presentation of graph-labelled edges.
